\subsection{Hướng phát triển tương lai}

Mặc dù đã đạt được những kết quả khả quan trong giai đoạn hiện tại, hệ thống vẫn còn nhiều tiềm năng để mở rộng và nâng cấp. Nhóm đề xuất các hướng phát triển tiếp theo nhằm đưa dự án trở thành một giải pháp thương mại điện tử toàn diện hơn:

\subsubsection{Nâng cao năng lực AI Agent}
\begin{itemize}
    \item \textbf{Tương tác đa phương thức (Multimodal):} Phát triển khả năng nhận diện hình ảnh, cho phép khách hàng tìm kiếm sản phẩm bằng ảnh chụp hoặc tương tác qua giọng nói.
    \item \textbf{Cá nhân hóa sâu (Hyper-personalization):} Tích hợp các thuật toán học máy nâng cao để phân tích sâu hơn hành vi người dùng, từ đó đưa ra các gợi ý sản phẩm chính xác tuyệt đối theo sở thích và thói quen tiêu dùng.
    \item \textbf{Mở rộng Agentic Workflow:} Nâng cấp AI Agent để có thể thực hiện các hành động thay cho người dùng một cách an toàn, ví dụ như tự động áp dụng mã giảm giá tốt nhất hoặc thực hiện đổi trả hàng tự động qua chat.
\end{itemize}

\subsubsection{Mở rộng quy mô hệ thống}
\begin{itemize}
    \item \textbf{Hỗ trợ đa nền tảng:} Phát triển ứng dụng di động bản địa (Native Mobile App) sử dụng React Native để đồng bộ trải nghiệm với phiên bản Web, tận dụng các tính năng phần cứng của điện thoại.
    \item \textbf{Hệ thống phân tích thời gian thực (Real-time Analytics):} Tích hợp các công cụ như Apache Flink hoặc ClickHouse để phân tích dữ liệu dòng (stream processing), cung cấp báo cáo kinh doanh tức thì cho quản trị viên.
    \item \textbf{Tự động hóa hạ tầng (Infrastructure as Code):} Sử dụng Terraform hoặc Pulumi để tự động hóa hoàn toàn việc thiết lập hạ tầng Cloud, giúp hệ thống có khả năng triển khai nhanh chóng trên nhiều vùng địa lý (multi-region).
\end{itemize}

\subsubsection{Bổ sung tính năng nghiệp vụ}
\begin{itemize}
    \item \textbf{Hệ thống khách hàng thân thiết (Loyalty System):} Phát triển cơ chế tích điểm, phân hạng thành viên và các chương trình ưu đãi tự động dựa trên mức độ gắn bó của khách hàng.
    \item \textbf{Tích hợp đa cổng thanh toán:} Mở rộng kết nối với các ví điện tử phổ biến và hỗ trợ thanh toán quốc tế để đa dạng hóa lựa chọn cho người dùng.
    \item \textbf{Quản lý chuỗi cung ứng nâng cao:} Tích hợp thêm các module dự báo nhu cầu hàng tồn kho dựa trên AI để tối ưu hóa chi phí vận hành cho doanh nghiệp.
\end{itemize}

\vspace{10pt}
Với những hướng phát triển này, nhóm tin rằng hệ thống có thể đáp ứng được các nhu cầu khắt khe của thị trường bán lẻ hiện đại, mang lại giá trị thực tiễn cao cho cả người dùng cuối và đơn vị kinh doanh.
